\section{Other options and their place in the investigation}
\label{sec:alternatives}

Analytic marginalization should come before additional approximation when
the model permits it. Section~\ref{sec:model-mixtures} gives the one-step
Gaussian example. A conditionally linear-Gaussian state block can similarly
be integrated rather than sampled, provided its conditional model is retained
exactly. A Gaussian approximation to a nonlinear block is a different choice
and needs its own error assessment. Neither a good fit to state observations
nor an appealing proposal density establishes that the downstream score is
accurate.

More particles, paired random designs, and randomized quasi-Monte Carlo are
important equal-compute alternatives. They can improve integration without
changing the scientific target, but they cannot restore a missing derivative
or correct a filtering approximation with persistent bias. The comparison
must distinguish a better estimator from simply spending more work on it.
Sparse mixtures, low-rank density approximations, and approximate backward
sums likewise require separate bias and cost measurements; they are not
algebraically equivalent to the full mixture merely because they produce
finite values.

Bandwidth or particle-count extrapolation is another possible bias-reduction
strategy, but it requires a justified expansion. A presumed $h^2$ smoothing
bias or $1/N$ score bias cannot be used across every nonsmooth, constrained,
or resampling model without evidence. Randomized telescoping is more general.
For coupled levels $S_l$ and an independent random truncation $L$,
\begin{equation}
 \widehat S_{\rm tel}
 =S_0+\sum_{l=1}^{L}\frac{S_l-S_{l-1}}{\Pr(L\ge l)},
 \qquad
 \mathbb E\widehat S_{\rm tel}=\lim_{l\to\infty}\mathbb E S_l,
 \label{eq:randomized-telescope}
\end{equation}
when absolute integrability permits exchanging sum and expectation. The
expectation equality follows because each retained increment has inclusion
probability $\Pr(L\ge l)$. The limit must first be shown to equal the model
score. Finite expected work and finite variance impose additional conditions
on the coupling and truncation tail; the displayed equality supplies neither.
This makes literal unbiasedness a distinct research problem rather than a
promised property of the next mixture prototype.

The original proposal-force use remains valid in a different role. A
deterministic position-only field may generate a reversible,
volume-preserving kick--drift--kick proposal, with acceptance computed from
the declared endpoint Hamiltonian. Each kick and drift is a shear with unit
Jacobian determinant. Reversing the step size reverses the sequence, and a
momentum flip turns the symmetric sequence into an involution. The ordinary
Metropolis ratio therefore corrects this proposal for the endpoint density;
the force need not be its exact gradient. If the endpoint uses a fixed
finite-particle likelihood, it preserves that approximate target, not the
exact state-space posterior. Exact pseudo-marginal inference instead needs a
valid nonnegative unbiased likelihood estimate and the auxiliary randomness
in the extended target. A new random score at each leapfrog step is not
automatically a reversible proposal.

The BayesFilter HMC interface already distinguishes an exact score from a
frozen proposal field. A future raw-coordinate field must use the typed
endpoint-corrected binding with \texttt{tune\_hmc\_kernel}; the
fixed-transport tuner is for a genuine frozen coordinate transport with its
Jacobian-corrected target. This document does not change an HMC consumer or
claim posterior readiness. The interface and capability registry were read
to keep the proposed downstream use consistent with the actual repository.

The practical order is therefore to check model identity and complete
initialization, establish a corrected mixture/Fisher reference, compare
selective gradient estimators where needed, and then investigate scalable
ancestor approximations or learned proposals. Control variates, analytic
integration, and increased integration effort should be evaluated alongside
that sequence. Further structural DSGE implementation remains deferred while
the existing repository models are checked; the support derivations below
are retained because they explain which model extensions will be valid.
